General form:  g(x) = a \cdot f\!\left(\frac{1}{b}(x - h)\right) + k

\begin{array}{ll}
|a| > 1 & vertical stretch by factor  |a| \\
0 < |a| < 1 & vertical compression \\
a < 0 & reflect over  x-axis \\
b > 1 & horizontal stretch by factor  b \\
0 < b < 1 & horizontal compression \\
h > 0 & shift right by  h \\
k > 0 & shift up by  k \\
\end{array}

Key point  (x_0,\, y_0)  on  f  maps to  (x_1,\, y_1)  on  g:
x_1 = b \cdot x_0 + h, \quad y_1 = a \cdot y_0 + k

Example (2026 \#8):  g(x) = -3f\!\left(\tfrac{1}{2}(x-6)\right) + 1

f  has maximum at  (-2,\, 5). Find max of  g.
x_1 = 2(-2) + 6 = 2, \quad y_1 = -3(5) + 1 = -14

Since  a = -3 < 0, max of  f  maps to a min of  g.
So  g  has a minimum at  (2, -14).

f  has minimum at  (4, -6):
x_1 = 2(4) + 6 = 14, \quad y_1 = -3(-6) + 1 = 19
g  has maximum at  (14,\, 19), so  w = 19

Domain shift:  f(x-h)+k \rightarrow domain shifts right by  h, range shifts up by  k

If  f  has domain  [-2,5], range  [-1,7]:
f(x-2)+1 \rightarrow domain  [0,7], range  [0,8]
a+b+c+d = 0+7+0+8 = 15

Vertex of  y = ax^2 + bx + c:
x_v = \frac{-b}{2a} \qquad y_v = c - \frac{b^2}{4a}

Maximum or minimum of quadratic:
y_{\max/\min} = \frac{4ac - b^2}{4a}
Maximum when  a < 0, minimum when  a > 0.
